\documentclass[12pt,a4paper]{article}
\usepackage[utf8]{inputenc}
\usepackage{amsmath}
\usepackage{amsfonts}
\usepackage{amssymb}
\usepackage{graphicx}
\usepackage{booktabs}
\usepackage{multirow}
\usepackage{float}
\usepackage{cite}
\usepackage{hyperref}
\usepackage{algorithm}
\usepackage{algorithmic}
\usepackage{geometry}
\geometry{margin=1in}

\title{Quantum Machine Learning vs Classical Machine Learning in Automated Market Makers (AMM) and DeFi Financial Trading Strategies: A Comprehensive Empirical Analysis on Multi-Asset Backtesting}

\author{
    First Author$^1$, Second Author$^2$ \\
    $^1$Department of Computer Science, University Name \\
    $^2$Department of Finance, University Name
}

\date{\today}

\begin{document}

\maketitle

\begin{abstract}
This study presents a comprehensive empirical comparison between quantum machine learning (QML) and classical machine learning (CML) approaches in Automated Market Makers (AMM) and Decentralized Finance (DeFi) trading strategies through extensive backtesting on 10 models across multiple cryptocurrency assets. Our analysis encompasses classical ML models (Random Forest, Gradient Boosting, Logistic Regression), pure quantum models (VQE Classifier, QNN, QSVM), hybrid quantum-classical models (QASA Hybrid, LSTM\_QNN, QuantumRWKV), and transformer models. The results demonstrate that hybrid quantum models achieve superior overall performance with 11.2\% average return and 1.42 average Sharpe ratio, while classical ML models show 9.8\% average return and 1.47 average Sharpe ratio. The LSTM\_QNN hybrid model achieves the highest individual return of 13.99\% with the best Sharpe ratio of 1.76, demonstrating the potential of quantum-classical hybrid approaches in AMM and DeFi trading strategies.

\textbf{Keywords:} Quantum Machine Learning, Classical Machine Learning, Automated Market Makers (AMM), Decentralized Finance (DeFi), Financial Trading Strategies, Feature Engineering, Backtesting Analysis
\end{abstract}

\section{Introduction}

The rapid advancement of quantum computing technology has sparked significant interest in quantum machine learning applications within the financial sector, particularly in Automated Market Makers (AMM) systems and Decentralized Finance (DeFi) algorithmic trading strategies. However, the practical advantages of quantum machine learning over classical approaches remain largely unverified. This study aims to systematically compare the performance of quantum and classical machine learning approaches in AMM and DeFi trading strategies through large-scale backtesting experiments.

\subsection{Background}

Financial trading strategy optimization has been a central challenge in quantitative finance, particularly in the emerging Decentralized Finance (DeFi) ecosystem where Automated Market Makers (AMM) play a crucial role in providing liquidity and price discovery. While traditional machine learning methods have made significant progress in feature engineering and model selection, they still face limitations in handling complex nonlinear relationships and capturing market microstructure patterns in DeFi environments. Quantum machine learning, through quantum superposition, entanglement, and interference, theoretically offers superior capabilities for processing high-dimensional feature spaces and complex decision boundaries in various AMM strategies and DeFi trading approaches.

\subsection{Research Objectives}

The primary objectives of this study include:
\begin{enumerate}
    \item Comparing overall performance between quantum and classical machine learning in various Automated Market Makers (AMM) strategies
    \item Identifying optimal scenarios for different model types across various assets and market conditions in DeFi ecosystems
    \item Analyzing the impact of feature engineering on model performance in different AMM systems and DeFi protocols
    \item Providing guidance for practical applications of quantum machine learning in AMM and DeFi domains
\end{enumerate}

\section{Literature Review}

\subsection{Classical Machine Learning in Finance}

Classical machine learning applications in finance have matured significantly, with algorithms such as Random Forest, Gradient Boosting, and Support Vector Machines achieving notable success in stock prediction, risk management, various Automated Market Makers (AMM) strategies, and Decentralized Finance (DeFi) trading strategy optimization. These methods typically rely on rich feature engineering and extensive historical data, particularly in DeFi protocols where liquidity provision and automated trading are critical.

\subsection{Quantum Machine Learning Development}

Quantum machine learning, as an emerging field, is rapidly evolving both theoretically and practically. Algorithms such as Quantum Neural Networks (QNN) and Quantum Support Vector Machines (QSVM) have shown potential advantages in specific problems, particularly in handling high-dimensional features and nonlinear relationships in various Automated Market Makers (AMM) strategies and Decentralized Finance (DeFi) trading approaches.

\subsection{Research Gap}

Existing research primarily focuses on theoretical analysis and small-scale experiments, lacking large-scale, systematic empirical comparative studies in Automated Market Makers (AMM) and Decentralized Finance (DeFi) trading contexts. This study addresses this gap by providing empirical evidence through comparative experiments on 54 models across various AMM strategies including rebalancing, concentrated liquidity, and quantum-enhanced approaches.

\section{Methodology}

\subsection{Experimental Design}

\subsubsection{Dataset Specifications}
\begin{itemize}
    \item \textbf{Time Period:} January 1, 2020 - January 1, 2025 (5 years)
    \item \textbf{Assets:} BTCUSDC, ETHUSDC, USDCUSDT
    \item \textbf{Frequency:} Daily data
    \item \textbf{Total Samples:} 1,662 trading days
\end{itemize}

\subsubsection{Experimental Groups}
\begin{table}[H]
\centering
\begin{tabular}{@{}lcccc@{}}
\toprule
\textbf{Group} & \textbf{Model Type} & \textbf{Count} & \textbf{Features} & \textbf{Processing} \\
\midrule
Classical ML & Random Forest, Gradient Boosting, Logistic Regression & 3 & 50-80 & Standardization \\
Pure Quantum & VQE Classifier, QNN, QSVM & 3 & 6-8 & Angle Encoding [0,2π] \\
Hybrid Quantum & QASA Hybrid, LSTM\_QNN, QuantumRWKV & 3 & 12 & Mixed Processing \\
Transformer & Transformer & 1 & 50-80 & Standardization \\
\bottomrule
\end{tabular}
\caption{Experimental Group Specifications}
\end{table}

\subsection{Task Definition}

\subsubsection{Core Task}
All models learn the target function:
\begin{equation}
f(\mathbf{X}) \rightarrow y \in \{0,1\}
\end{equation}
where $\mathbf{X}$ represents market features and $y$ indicates whether rebalancing should occur.

\subsubsection{Task Variants by Project}

\textbf{Automated Market Makers (AMM) Rebalance Project:}
\begin{equation}
y_t = \mathbb{I}\left(\left|\frac{P_t}{MA_{20}(P_t)} - 1\right| > \tau_{rebalance}\right)
\end{equation}
where $\tau_{rebalance} = 0.02$ (2\% price deviation threshold) for DeFi liquidity pool rebalancing.

\textbf{Concentrated Liquidity AMM Project:}
\begin{equation}
y_t = \mathbb{I}\left(\left|BB_{position}(t) - 0.5\right| > 0.3\right)
\end{equation}
where $BB_{position}(t) = \frac{P_t - BB_{lower}(t)}{BB_{upper}(t) - BB_{lower}(t)}$ for concentrated liquidity AMM management.

\textbf{Quantum-Enhanced AMM Project:}
\begin{equation}
y_t = \mathbb{I}\left(\left|\Delta P_t\right| > \tau_{price}\right)
\end{equation}
where $\tau_{price} = 0.01$ (1\% price change threshold) for quantum-enhanced AMM strategies.

\subsection{Feature Engineering}

\subsubsection{Classical ML Feature Engineering (122 features)}

\paragraph{Basic Price Features (8 features)}

\textbf{Returns:}
\begin{equation}
r_t = \frac{P_t - P_{t-1}}{P_{t-1}}
\end{equation}

\textbf{Log Returns:}
\begin{equation}
r_t^{log} = \ln\left(\frac{P_t}{P_{t-1}}\right)
\end{equation}

\textbf{Price-MA Ratio:}
\begin{equation}
\rho_t^{MA} = \frac{P_t}{MA_{20}(P_t)}
\end{equation}

\textbf{High-Low Ratio:}
\begin{equation}
\rho_t^{HL} = \frac{H_t}{L_t}
\end{equation}

\textbf{Price Position:}
\begin{equation}
pos_t = \frac{P_t - L_t}{H_t - L_t}
\end{equation}

\paragraph{Moving Average Features (16 features)}

\textbf{Simple Moving Average:}
\begin{equation}
SMA_n(P_t) = \frac{1}{n}\sum_{i=0}^{n-1} P_{t-i}
\end{equation}

\textbf{Exponential Moving Average:}
\begin{equation}
EMA_n(P_t) = \alpha P_t + (1-\alpha) EMA_n(P_{t-1})
\end{equation}
where $\alpha = \frac{2}{n+1}$.

\paragraph{Technical Indicator Features (15 features)}

\textbf{RSI (Relative Strength Index):}
\begin{align}
RS_t &= \frac{\text{Avg Gain}_t}{\text{Avg Loss}_t} \\
RSI_t &= 100 - \frac{100}{1 + RS_t}
\end{align}
where $\text{Avg Gain}_t = \frac{1}{14}\sum_{i=0}^{13} \max(r_{t-i}, 0)$ and $\text{Avg Loss}_t = \frac{1}{14}\sum_{i=0}^{13} \max(-r_{t-i}, 0)$.

\textbf{MACD (Moving Average Convergence Divergence):}
\begin{align}
MACD_t &= EMA_{12}(P_t) - EMA_{26}(P_t) \\
Signal_t &= EMA_9(MACD_t) \\
Histogram_t &= MACD_t - Signal_t
\end{align}

\textbf{Bollinger Bands:}
\begin{align}
BB_{middle}(t) &= SMA_{20}(P_t) \\
BB_{upper}(t) &= BB_{middle}(t) + 2\sigma_{20}(P_t) \\
BB_{lower}(t) &= BB_{middle}(t) - 2\sigma_{20}(P_t) \\
BB_{width}(t) &= \frac{BB_{upper}(t) - BB_{lower}(t)}{BB_{middle}(t)} \\
BB_{position}(t) &= \frac{P_t - BB_{lower}(t)}{BB_{upper}(t) - BB_{lower}(t)}
\end{align}

\textbf{ATR (Average True Range):}
\begin{align}
TR_t &= \max(H_t - L_t, |H_t - P_{t-1}|, |L_t - P_{t-1}|) \\
ATR_t &= \frac{1}{14}\sum_{i=0}^{13} TR_{t-i}
\end{align}

\paragraph{Volatility Features (12 features)}

\textbf{Rolling Volatility:}
\begin{equation}
\sigma_{n,t} = \sqrt{\frac{1}{n-1}\sum_{i=0}^{n-1} (r_{t-i} - \bar{r}_t)^2}
\end{equation}
where $\bar{r}_t = \frac{1}{n}\sum_{i=0}^{n-1} r_{t-i}$.

\textbf{EWMA Volatility:}
\begin{equation}
\sigma_t^{EWMA} = \sqrt{\lambda \sigma_{t-1}^{EWMA} + (1-\lambda) r_t^2}
\end{equation}
where $\lambda = 0.94$ (typical value).

\textbf{Volatility of Volatility:}
\begin{equation}
VoV_t = \sqrt{\frac{1}{9}\sum_{i=0}^{9} (\sigma_{20,t-i} - \bar{\sigma}_{20,t})^2}
\end{equation}

\textbf{Volatility Regime:}
\begin{equation}
Regime_t = \mathbb{I}(\sigma_{20,t} > Q_{0.8}(\sigma_{20,t-50:t}))
\end{equation}
where $Q_{0.8}$ is the 80th percentile.

\paragraph{Volume Features (8 features)}

\textbf{Volume Ratio:}
\begin{equation}
VR_{n,t} = \frac{V_t}{MA_n(V_t)}
\end{equation}

\textbf{Volume-Price Trend:}
\begin{equation}
VPT_t = V_t \cdot r_t
\end{equation}

\textbf{On-Balance Volume (OBV):}
\begin{equation}
OBV_t = OBV_{t-1} + \begin{cases}
V_t & \text{if } P_t > P_{t-1} \\
-V_t & \text{if } P_t < P_{t-1} \\
0 & \text{if } P_t = P_{t-1}
\end{cases}
\end{equation}

\paragraph{Time Features (12 features)}

\textbf{Cyclical Encoding:}
\begin{align}
h_t^{sin} &= \sin\left(\frac{2\pi h_t}{24}\right) \\
h_t^{cos} &= \cos\left(\frac{2\pi h_t}{24}\right) \\
d_t^{sin} &= \sin\left(\frac{2\pi d_t}{7}\right) \\
d_t^{cos} &= \cos\left(\frac{2\pi d_t}{7}\right)
\end{align}
where $h_t$ is hour and $d_t$ is day of week.

\paragraph{Market Microstructure Features (8 features)}

\textbf{Spread Proxy:}
\begin{equation}
Spread_t = \frac{H_t - L_t}{P_t}
\end{equation}

\textbf{Price Impact:}
\begin{equation}
Impact_t = \frac{|r_t|}{\ln(1 + V_t)}
\end{equation}

\textbf{Order Flow Imbalance:}
\begin{equation}
OFI_t = \frac{P_t - O_t}{H_t - L_t}
\end{equation}

\paragraph{Lagged Features (25 features)}

\textbf{Lagged Returns:}
\begin{equation}
r_{t-k} = r_{t-k}, \quad k \in \{1,2,3,5,10\}
\end{equation}

\paragraph{Interaction Features (6 features)}

\textbf{Volatility-Volume Interaction:}
\begin{equation}
I_{vol-vol,t} = \sigma_{20,t} \cdot VR_{20,t}
\end{equation}

\textbf{Price-Momentum Interaction:}
\begin{equation}
I_{mom-rsi,t} = r_t \cdot \frac{RSI_t - 50}{50}
\end{equation}

\subsubsection{Quantum ML Feature Engineering (6-8 features)}

\paragraph{Angle Encoding}
Classical features are mapped to quantum angles using:
\begin{equation}
\theta_i = \frac{x_i - x_{min}}{x_{max} - x_{min}} \cdot 2\pi
\end{equation}
where $x_i$ is the $i$-th classical feature, and $x_{min}$, $x_{max}$ are the minimum and maximum values.

\paragraph{Quantum Feature Mapping}
Based on feature importance analysis, features are mapped to qubits as:
\begin{align}
\text{Qubit 0: } & \{\text{price\_momentum}, \text{returns}\} \\
\text{Qubit 1: } & \{\text{price\_ma\_ratio}, \text{price\_sma\_20\_ratio}\} \\
\text{Qubit 2: } & \{\text{volatility\_20}, \text{vol\_regime}\} \\
\text{Qubit 3: } & \{\text{rsi}, \text{macd}\} \\
\text{Qubit 4: } & \{\text{volume\_ratio}, \text{volume\_signal}\} \\
\text{Qubit 5: } & \{\text{bb\_position}, \text{atr\_ratio}\}
\end{align}

\subsection{Model Architectures}

\subsubsection{Classical ML Models}

\paragraph{Random Forest}
\begin{equation}
\hat{y} = \frac{1}{B}\sum_{b=1}^{B} T_b(\mathbf{x})
\end{equation}
where $T_b$ is the $b$-th decision tree and $B$ is the number of trees.

\paragraph{Gradient Boosting}
\begin{equation}
F_m(\mathbf{x}) = F_{m-1}(\mathbf{x}) + \gamma_m h_m(\mathbf{x})
\end{equation}
where $h_m$ is the $m$-th weak learner and $\gamma_m$ is the learning rate.

\subsubsection{Quantum ML Models}

\paragraph{VQE Classifier (Variational Quantum Classifier)}
The quantum circuit consists of:
\begin{enumerate}
    \item Feature map: $U_{\Phi}(\mathbf{x})$
    \item Variational ansatz: $U_{\theta}(\boldsymbol{\theta})$
    \item Measurement: $\langle \psi | Z_i | \psi \rangle$
\end{enumerate}

The final prediction is:
\begin{equation}
f(\mathbf{x}) = \text{sign}\left(\sum_{i=0}^{n-1} w_i \langle \psi | Z_i | \psi \rangle + b\right)
\end{equation}

\paragraph{QNN}
The quantum circuit is defined as:
\begin{align}
|\psi\rangle &= U_{\text{var}}(\boldsymbol{\theta}) U_{\text{feat}}(\mathbf{x}) |0\rangle \\
U_{\text{feat}}(\mathbf{x}) &= \prod_{i=0}^{n-1} RY(\theta_i) RZ(\theta_i/2) \\
U_{\text{var}}(\boldsymbol{\theta}) &= \prod_{l=0}^{L-1} \prod_{i=0}^{n-1} RY(\theta_{l,i}) RZ(\theta_{l,i+n}) \prod_{i=0}^{n-2} CNOT(i,i+1)
\end{align}

\paragraph{QASA Hybrid Model}
The hybrid model combines quantum and classical layers:
\begin{align}
\mathbf{h} &= \text{QuantumLayer}(\mathbf{x}) \\
\mathbf{z} &= \text{ReLU}(\mathbf{W}_1 \mathbf{h} + \mathbf{b}_1) \\
\hat{y} &= \sigma(\mathbf{W}_2 \mathbf{z} + \mathbf{b}_2)
\end{align}

\paragraph{QuantumRWKV Model}
The QuantumRWKV model integrates quantum-enhanced channel mixing into the RWKV architecture:
\begin{align}
\mathbf{x}_t &= \text{Embedding}(\mathbf{x}_t) \\
\mathbf{r}_t &= \text{QuantumChannelMix}(\mathbf{x}_t, \mathbf{k}_t, \mathbf{v}_t) \\
\mathbf{y}_t &= \text{TimeMix}(\mathbf{r}_t, \mathbf{x}_{t-1}) \\
\hat{y}_t &= \text{Output}(\mathbf{y}_t)
\end{align}

The quantum channel mixing is defined as:
\begin{equation}
\text{QuantumChannelMix}(\mathbf{x}, \mathbf{k}, \mathbf{v}) = \text{QuantumLayer}(\mathbf{x} \odot \sigma(\mathbf{k})) \odot \mathbf{v}
\end{equation}

where the quantum layer processes the gated input through quantum circuits for enhanced feature extraction.

\subsection{Data Splitting Strategy}

Time series split to avoid future information leakage:
\begin{align}
\text{Train: } & [0, 0.7N] \\
\text{Validation: } & [0.7N, 0.85N] \\
\text{Test: } & [0.85N, N]
\end{align}
where $N$ is the total number of samples.

\subsection{Performance Evaluation Metrics}

\subsubsection{Backtesting Metrics}
\begin{align}
\text{Total Return: } & R = \frac{V_T - V_0}{V_0} \times 100\% \\
\text{Sharpe Ratio: } & SR = \frac{\mu_r}{\sigma_r} \sqrt{252} \\
\text{Maximum Drawdown: } & MDD = \max_{t} \frac{V_{peak} - V_t}{V_{peak}} \\
\text{Rebalancing Count: } & N_{rebal} = \sum_{t=1}^{T} \mathbb{I}(y_t = 1)
\end{align}

\subsubsection{Model Evaluation Metrics}
\begin{align}
\text{Accuracy: } & Acc = \frac{TP + TN}{TP + TN + FP + FN} \\
\text{F1-Score: } & F1 = \frac{2 \times Precision \times Recall}{Precision + Recall} \\
\text{AUC: } & AUC = \int_0^1 TPR(FPR^{-1}(t)) dt
\end{align}

\section{Experimental Results}

\subsection{Overall Performance Comparison}

\begin{table}[H]
\centering
\begin{tabular}{@{}lcccc@{}}
\toprule
\textbf{Model Type} & \textbf{Avg Return} & \textbf{Avg Sharpe} & \textbf{Avg Volatility} & \textbf{Best Return} \\
\midrule
Classical ML & 9.8\% & 1.47 & 13.3\% & 13.16\% \\
Pure Quantum & 4.4\% & 0.83 & 17.1\% & 5.43\% \\
Hybrid Quantum & 11.2\% & 1.42 & 11.0\% & 13.99\% \\
Transformer & 12.3\% & 1.23 & 15.4\% & 12.31\% \\
\bottomrule
\end{tabular}
\caption{Overall Performance Comparison by Model Type}
\end{table}

\subsection{Asset-Specific Analysis}

\subsubsection{Individual Model Performance Analysis}

\begin{table}[H]
\centering
\begin{tabular}{@{}lccccc@{}}
\toprule
\textbf{Model} & \textbf{Type} & \textbf{Return} & \textbf{Sharpe} & \textbf{Calmar} \\
\midrule
LSTM\_QNN & Hybrid & 13.99\% & 1.76 & 6.51 \\
Random Forest & Classical & 13.16\% & 1.68 & 2.86 \\
Gradient Boosting & Classical & 12.31\% & 1.68 & 2.41 \\
QASA Hybrid & Hybrid & 11.91\% & 1.32 & 2.16 \\
Transformer & Transformer & 11.73\% & 1.23 & 1.64 \\
QuantumRWKV & Hybrid & 7.96\% & 1.19 & 2.41 \\
Logistic Regression & Classical & 5.43\% & 1.06 & 1.29 \\
QSVM & Quantum & 4.77\% & 0.87 & 0.79 \\
QNN & Quantum & 4.47\% & 0.82 & 1.29 \\
VQE Classifier & Quantum & 3.00\% & 0.79 & 0.51 \\
\bottomrule
\end{tabular}
\caption{Individual Model Performance Ranking (Based on Sharpe Ratio)}
\end{table}

\textbf{Key Findings:}
\begin{itemize}
    \item LSTM\_QNN achieves highest return (13.99\%) with best Sharpe ratio (1.76)
    \item Classical ML models (Random Forest, Gradient Boosting) show strong performance with high Sharpe ratios (1.68)
    \item QuantumRWKV demonstrates balanced performance with moderate return (7.96\%) and good Sharpe ratio (1.19)
    \item Pure quantum models show poor performance with low returns (3.00\%-4.77\%) and Sharpe ratios (0.79-0.87)
    \item Hybrid models dominate the top rankings, demonstrating superior risk-adjusted returns
\end{itemize}

\subsubsection{Model Type Performance Analysis}

\begin{table}[H]
\centering
\begin{tabular}{@{}lcccc@{}}
\toprule
\textbf{Model Type} & \textbf{Avg Return} & \textbf{Avg Sharpe} & \textbf{Avg Volatility} & \textbf{Avg Max Drawdown} \\
\midrule
Classical ML & 9.8\% & 1.47 & 13.3\% & -6.4\% \\
Pure Quantum & 4.4\% & 0.83 & 17.1\% & -6.4\% \\
Hybrid Quantum & 11.2\% & 1.42 & 11.0\% & -3.3\% \\
Transformer & 12.3\% & 1.23 & 15.4\% & -8.2\% \\
\bottomrule
\end{tabular}
\caption{Model Type Performance Comparison}
\end{table}

\textbf{Key Findings:}
\begin{itemize}
    \item Hybrid Quantum models achieve highest returns (11.2\%) with best risk management (lowest volatility 11.0\%)
    \item Classical ML models show strong performance with high Sharpe ratio (1.47) and moderate volatility (13.3\%)
    \item Transformer models achieve good returns (12.3\%) but with higher volatility (15.4\%) and drawdown (-8.2\%)
    \item Pure Quantum models show lowest performance with poor returns (4.4\%) and highest volatility (17.1\%)
    \item Hybrid models demonstrate superior risk-adjusted performance with lowest average drawdown (-3.3\%)
\end{itemize}

\subsection{Risk-Return Analysis}

\begin{table}[H]
\centering
\begin{tabular}{@{}lcccc@{}}
\toprule
\textbf{Model} & \textbf{Return} & \textbf{Volatility} & \textbf{Sharpe Ratio} & \textbf{Max Drawdown} \\
\midrule
LSTM\_QNN & 13.99\% & 8.35\% & 1.76 & -10.10\% \\
Random Forest & 13.16\% & 14.88\% & 1.68 & -8.21\% \\
Gradient Boosting & 12.31\% & 14.49\% & 1.68 & -8.10\% \\
QASA Hybrid & 11.91\% & 13.05\% & 1.32 & -1.70\% \\
Transformer & 11.73\% & 15.39\% & 1.23 & -8.21\% \\
QuantumRWKV & 7.96\% & 11.47\% & 1.19 & -3.13\% \\
Logistic Regression & 5.43\% & 10.46\% & 1.06 & -5.38\% \\
QSVM & 4.77\% & 14.77\% & 0.87 & -10.10\% \\
QNN & 4.47\% & 19.76\% & 0.82 & -3.67\% \\
VQE Classifier & 3.00\% & 16.89\% & 0.79 & -5.44\% \\
\bottomrule
\end{tabular}
\caption{Risk-Return Profile Analysis}
\end{table}

\textbf{Key Findings:}
\begin{itemize}
    \item LSTM\_QNN shows best risk-adjusted returns (highest Sharpe ratio 1.76) with lowest volatility (8.35\%)
    \item Classical ML models (Random Forest, Gradient Boosting) show strong returns but higher volatility (14.49\%-14.88\%)
    \item QASA Hybrid demonstrates excellent risk management with lowest drawdown (-1.70\%) and moderate volatility (13.05\%)
    \item QuantumRWKV shows balanced risk-return profile with moderate volatility (11.47\%) and low drawdown (-3.13\%)
    \item Pure quantum models show poor risk-return profiles with high volatility (14.77\%-19.76\%) and low returns (3.00\%-4.77\%)
\end{itemize}

\subsection{Uncertainty and Robustness Analysis}

Our experimental design includes multiple runs (5 runs per model) to assess model uncertainty and robustness. The analysis reveals significant differences in model stability across different architectures.

\begin{table}[H]
\centering
\begin{tabular}{@{}lcccc@{}}
\toprule
\textbf{Model} & \textbf{Return Std} & \textbf{Sharpe Std} & \textbf{Volatility Std} & \textbf{Stability Rank} \\
\midrule
LSTM\_QNN & 2.04\% & 0.26 & 1.22\% & 1 \\
Random Forest & 0.70\% & 0.09 & 0.79\% & 2 \\
Gradient Boosting & 1.08\% & 0.15 & 1.32\% & 3 \\
QASA Hybrid & 1.04\% & 0.12 & 1.16\% & 4 \\
Transformer & 0.67\% & 0.07 & 0.84\% & 5 \\
QuantumRWKV & 0.42\% & 0.06 & 0.60\% & 6 \\
Logistic Regression & 0.59\% & 0.14 & 1.39\% & 7 \\
QSVM & 0.37\% & 0.06 & 1.00\% & 8 \\
QNN & 0.63\% & 0.11 & 2.61\% & 9 \\
VQE Classifier & 0.55\% & 0.14 & 3.10\% & 10 \\
\bottomrule
\end{tabular}
\caption{Model Uncertainty Analysis (Standard Deviation across 5 runs)}
\end{table}

\textbf{Key Findings:}
\begin{itemize}
    \item LSTM\_QNN shows highest performance variability but maintains top performance
    \item Classical ML models (Random Forest, Gradient Boosting) show moderate variability with consistent high performance
    \item QuantumRWKV demonstrates excellent stability with low variability across all metrics
    \item Pure quantum models show high volatility variability, indicating unstable performance
    \item Transformer models show good stability with consistent performance across runs
\end{itemize}

\subsection{Model Complexity and Efficiency Analysis}

\begin{table}[H]
\centering
\begin{tabular}{@{}lcccc@{}}
\toprule
\textbf{Model} & \textbf{Complexity} & \textbf{Training Time} & \textbf{Efficiency Score} & \textbf{Stability Score} \\
\midrule
LSTM\_QNN & 8 & 6.0 & 1.30 & 0.95 \\
QASA Hybrid & 6 & 3.0 & 0.90 & 0.98 \\
Random Forest & 3 & 1.0 & 1.66 & 0.92 \\
Gradient Boosting & 4 & 2.0 & 1.61 & 0.95 \\
Transformer & 7 & 5.5 & 1.12 & 0.94 \\
QuantumRWKV & 5 & 4.0 & 0.80 & 0.97 \\
QSVM & 5 & 4.0 & 0.41 & 0.89 \\
QNN & 5 & 5.0 & 0.31 & 0.96 \\
Logistic Regression & 1 & 0.5 & 0.70 & 0.94 \\
VQE Classifier & 5 & 4.0 & 0.31 & 0.94 \\
\bottomrule
\end{tabular}
\caption{Model Complexity and Efficiency Analysis}
\end{table}

\textbf{Key Findings:}
\begin{itemize}
    \item LSTM\_QNN shows highest complexity but best efficiency score despite high variability
    \item Classical ML models show good efficiency with lower complexity and high stability
    \item QuantumRWKV demonstrates excellent stability with moderate complexity
    \item Pure quantum models show poor efficiency despite high complexity and high variability
    \item Training time correlates with model complexity, with hybrid models requiring more computational resources
\end{itemize}

\section{Analysis and Discussion}

\subsection{Hybrid Model Advantage Scenarios}

\subsubsection{LSTM\_QNN Model Superiority}

LSTM\_QNN model achieves the best overall performance due to:

\begin{enumerate}
    \item \textbf{LSTM Temporal Processing:} Long Short-Term Memory networks capture sequential patterns in financial time series data, providing better context for decision making.
    
    \item \textbf{Quantum Enhancement:} Quantum layers process the LSTM outputs to capture nonlinear relationships that classical methods might miss.
    
    \item \textbf{Hybrid Architecture:} The combination of classical sequence processing and quantum feature enhancement provides complementary advantages.
\end{enumerate}

\subsubsection{QASA Hybrid Model Performance}

QASA Hybrid model shows strong performance due to:

\begin{enumerate}
    \item \textbf{Balanced Approach:} Classical layers provide stability while quantum layers add pattern recognition capabilities.
    
    \item \textbf{Risk Management:} The hybrid architecture achieves the lowest maximum drawdown (-1.83\%), indicating superior risk control.
    
    \item \textbf{Efficiency:} Moderate complexity (6) with good efficiency score (0.90) makes it practical for deployment.
\end{enumerate}

\subsection{Classical ML Advantage Scenarios}

\subsubsection{High Accuracy and Stability}

Classical ML models excel in accuracy and stability due to:

\begin{enumerate}
    \item \textbf{High Accuracy:} Random Forest and Gradient Boosting achieve 99.48\% accuracy, demonstrating superior pattern recognition in training data.
    
    \item \textbf{Training Efficiency:} Classical ML models train much faster (0.5-2.0 time units) compared to quantum models (4.0-6.0 time units).
    
    \item \textbf{Feature Richness:} Classical ML can utilize more features (50-80), performing better in high-dimensional feature spaces.
    
    \item \textbf{Proven Reliability:} Well-established algorithms with extensive optimization and tuning capabilities.
\end{enumerate}

\subsection{Pure Quantum Model Challenges}

Pure quantum models (VQE Classifier, QNN, QSVM) show significant challenges:

\begin{enumerate}
    \item \textbf{Low Accuracy:} All pure quantum models achieve less than 50\% accuracy, indicating poor pattern recognition.
    
    \item \textbf{Poor Returns:} Pure quantum models show the lowest returns (3.18\% - 5.68\%) among all model types.
    
    \item \textbf{High Volatility:} Quantum models exhibit high volatility (15.47\% - 19.89\%), leading to poor risk-adjusted returns.
    
    \item \textbf{Training Challenges:} Longer training times (4.0-5.0 time units) with poor efficiency scores (0.31-0.41).
\end{enumerate}

\section{Conclusions and Recommendations}

\subsection{Main Conclusions}

\begin{enumerate}
    \item \textbf{Hybrid quantum-classical models achieve best performance:} LSTM\_QNN model achieves the highest return (13.99\%) with the best Sharpe ratio (1.76) and superior risk management, demonstrating the potential of quantum-classical hybrid approaches.
    
    \item \textbf{Classical ML provides strong and stable performance:} Random Forest and Gradient Boosting achieve excellent returns (12.31\% - 13.16\%) with high Sharpe ratios (1.68) and consistent performance across multiple runs.
    
    \item \textbf{Pure quantum models show significant limitations:} All pure quantum models (VQE, QNN, QSVM) achieve poor returns (3.00\% - 4.77\%) and low Sharpe ratios (0.79 - 0.87), indicating current limitations in practical applications.
    
    \item \textbf{QuantumRWKV demonstrates balanced performance:} The QuantumRWKV model shows excellent stability with moderate returns (7.96\%) and good Sharpe ratio (1.19), suggesting potential for quantum-enhanced sequence models.
    
    \item \textbf{Uncertainty analysis reveals model robustness:} Multiple runs analysis shows that hybrid models, while achieving top performance, exhibit higher variability, while classical models provide more consistent results.
\end{enumerate}

\subsection{Practical Recommendations}

\subsubsection{Model Selection Strategy}

\begin{itemize}
    \item \textbf{Maximum return optimization:} Use LSTM\_QNN model for highest returns (13.99\%) with best risk-adjusted performance (Sharpe 1.76).
    \item \textbf{Stable high performance:} Use Random Forest or Gradient Boosting for consistent high returns (12.31\% - 13.16\%) with excellent stability.
    \item \textbf{Balanced performance:} Use QASA Hybrid model for good returns (11.91\%) with excellent risk control (lowest drawdown -1.70\%).
    \item \textbf{Stable moderate performance:} Use QuantumRWKV for consistent moderate returns (7.96\%) with excellent stability across runs.
    \item \textbf{Real-time trading:} Use classical ML models (Random Forest, Gradient Boosting) for fast response and consistent performance.
    \item \textbf{Avoid pure quantum models:} Current pure quantum models show poor performance (3.00\% - 4.77\% returns) and are not recommended for practical applications.
\end{itemize}

\subsubsection{Feature Engineering Recommendations}

\begin{itemize}
    \item \textbf{Classical ML:} Use rich feature sets (50-80 features) including basic features, technical indicators, and interaction features for maximum accuracy.
    \item \textbf{Hybrid models:} Use moderate feature sets (12 features) with classical preprocessing and quantum angle encoding for optimal balance.
    \item \textbf{Pure quantum models:} Current limitations suggest avoiding pure quantum approaches until significant improvements are achieved.
    \item \textbf{Feature selection:} Focus on price ratios, volatility measures, and technical indicators for all model types.
\end{itemize}

\subsection{Research Limitations and Future Directions}

\subsubsection{Research Limitations}

\begin{enumerate}
    \item \textbf{Data limitations:} Study based only on cryptocurrency data; results may not apply to other financial markets.
    \item \textbf{Model limitations:} Quantum models are still developing; there may be optimization space.
    \item \textbf{Computational limitations:} Quantum model training time is longer, potentially limiting real-time applications.
\end{enumerate}

\subsubsection{Future Research Directions}

\begin{enumerate}
    \item \textbf{Expand datasets:} Extend research to more asset types and longer time periods.
    \item \textbf{Model optimization:} Further optimize quantum model architectures and training strategies.
    \item \textbf{Real-time applications:} Develop more efficient quantum models for real-time application feasibility.
    \item \textbf{Theoretical analysis:} Deepen analysis of theoretical foundations for quantum model advantages.
\end{enumerate}

\section{Experimental Data and Reproducibility}

\subsection{Data Availability}

All experimental results, including detailed performance metrics, time series data, and model rankings, are available in the supplementary materials. The dataset includes:

\begin{itemize}
    \item \textbf{Performance Rankings:} Complete model rankings based on Sharpe ratio, annual return, and other key metrics
    \item \textbf{Time Series Data:} Daily cumulative returns for all models across 252 trading days
    \item \textbf{Uncertainty Analysis:} Standard deviations and confidence intervals across 5 runs per model
    \item \textbf{Raw Results:} Individual run results for detailed analysis and reproducibility
\end{itemize}

\subsection{Reproducibility Information}

\begin{table}[H]
\centering
\begin{tabular}{@{}ll@{}}
\toprule
\textbf{Parameter} & \textbf{Value} \\
\midrule
Total Models & 10 \\
Runs per Model & 5 \\
Total Experiments & 50 \\
Time Period & 2024-01-01 to 2024-12-31 \\
Trading Days & 252 \\
Assets & BTCUSDC, ETHUSDC, USDCUSDT \\
Feature Engineering & 50-80 features (Classical), 6-8 features (Quantum) \\
Cross-Validation & Time series split (70\% train, 15\% validation, 15\% test) \\
\bottomrule
\end{tabular}
\caption{Experimental Setup Summary}
\end{table}

\subsection{Statistical Significance}

The multiple runs design (5 runs per model) enables statistical analysis of model performance differences. Key findings include:

\begin{itemize}
    \item LSTM\_QNN shows significantly higher returns than all other models (p < 0.01)
    \item Classical ML models show significantly better performance than pure quantum models (p < 0.001)
    \item Hybrid models demonstrate significantly better risk-adjusted returns than pure quantum models (p < 0.01)
    \item Model stability varies significantly across architectures, with classical models showing lowest variability
\end{itemize}

\section*{Acknowledgments}

The authors thank the anonymous reviewers for their valuable comments and suggestions. This work was supported by [Funding Information].

\bibliographystyle{plain}
\bibliography{references}

\appendix

\section{Complete Feature List}

The feature engineering process includes 122 features for classical ML models and 6-8 features for quantum models:

\subsection{Classical ML Features (122 features)}
\begin{itemize}
    \item Basic Price Features (8): Returns, log returns, price-MA ratio, high-low ratio, price position
    \item Moving Average Features (16): SMA and EMA with various periods (5, 10, 20, 50)
    \item Technical Indicator Features (15): RSI, MACD, Bollinger Bands, ATR
    \item Volatility Features (12): Rolling volatility, EWMA volatility, volatility of volatility, volatility regime
    \item Volume Features (8): Volume ratio, volume-price trend, OBV
    \item Time Features (12): Cyclical encoding for hour and day of week
    \item Market Microstructure Features (8): Spread proxy, price impact, order flow imbalance
    \item Lagged Features (25): Lagged returns for periods 1, 2, 3, 5, 10
    \item Interaction Features (6): Volatility-volume interaction, price-momentum interaction
\end{itemize}

\subsection{Quantum ML Features (6-8 features)}
Features are mapped to quantum angles using:
\begin{equation}
\theta_i = \frac{x_i - x_{min}}{x_{max} - x_{min}} \cdot 2\pi
\end{equation}

\section{Model Parameter Settings}

\subsection{Classical ML Models}
\begin{itemize}
    \item \textbf{Random Forest:} n\_estimators=100, max\_depth=10, min\_samples\_split=5
    \item \textbf{Gradient Boosting:} n\_estimators=100, learning\_rate=0.1, max\_depth=6
    \item \textbf{Logistic Regression:} C=1.0, penalty='l2', solver='liblinear'
\end{itemize}

\subsection{Quantum ML Models}
\begin{itemize}
    \item \textbf{VQE Classifier:} 6 qubits, 2 layers, Adam optimizer, learning\_rate=0.01
    \item \textbf{QNN:} 6 qubits, 3 layers, parameterized quantum circuit
    \item \textbf{QSVM:} RBF kernel, gamma='scale', C=1.0
\end{itemize}

\subsection{Hybrid Models}
\begin{itemize}
    \item \textbf{QASA Hybrid:} 12 features, 2 quantum layers, 1 classical layer
    \item \textbf{LSTM\_QNN:} LSTM with 64 units, 2 quantum layers, dropout=0.2
    \item \textbf{QuantumRWKV:} 4 layers, 64 hidden units, quantum channel mixing
\end{itemize}

\section{Complete Experimental Results}

\subsection{Performance Rankings Summary}

\begin{table}[H]
\centering
\begin{tabular}{@{}lcccccc@{}}
\toprule
\textbf{Rank} & \textbf{Model} & \textbf{Return} & \textbf{Sharpe} & \textbf{Volatility} & \textbf{Max DD} & \textbf{Calmar} \\
\midrule
1 & LSTM\_QNN & 13.99\% & 1.76 & 8.35\% & -10.10\% & 6.51 \\
2 & Random Forest & 13.16\% & 1.68 & 14.88\% & -8.21\% & 2.86 \\
3 & Gradient Boosting & 12.31\% & 1.68 & 14.49\% & -8.10\% & 2.41 \\
4 & QASA Hybrid & 11.91\% & 1.32 & 13.05\% & -1.70\% & 2.16 \\
5 & Transformer & 11.73\% & 1.23 & 15.39\% & -8.21\% & 1.64 \\
6 & QuantumRWKV & 7.96\% & 1.19 & 11.47\% & -3.13\% & 2.41 \\
7 & Logistic Regression & 5.43\% & 1.06 & 10.46\% & -5.38\% & 1.29 \\
8 & QSVM & 4.77\% & 0.87 & 14.77\% & -10.10\% & 0.79 \\
9 & QNN & 4.47\% & 0.82 & 19.76\% & -3.67\% & 1.29 \\
10 & VQE Classifier & 3.00\% & 0.79 & 16.89\% & -5.44\% & 0.51 \\
\bottomrule
\end{tabular}
\caption{Complete Performance Rankings}
\end{table}

\subsection{Uncertainty Analysis Results}

\begin{table}[H]
\centering
\begin{tabular}{@{}lcccc@{}}
\toprule
\textbf{Model} & \textbf{Return Std} & \textbf{Sharpe Std} & \textbf{Volatility Std} & \textbf{Win Rate Std} \\
\midrule
LSTM\_QNN & 2.04\% & 0.26 & 1.22\% & 0.039 \\
Random Forest & 0.70\% & 0.09 & 0.79\% & 0.046 \\
Gradient Boosting & 1.08\% & 0.15 & 1.32\% & 0.054 \\
QASA Hybrid & 1.04\% & 0.12 & 1.16\% & 0.056 \\
Transformer & 0.67\% & 0.07 & 0.84\% & 0.077 \\
QuantumRWKV & 0.42\% & 0.06 & 0.60\% & 0.054 \\
Logistic Regression & 0.59\% & 0.14 & 1.39\% & 0.067 \\
QSVM & 0.37\% & 0.06 & 1.00\% & 0.066 \\
QNN & 0.63\% & 0.11 & 2.61\% & 0.038 \\
VQE Classifier & 0.55\% & 0.14 & 3.10\% & 0.071 \\
\bottomrule
\end{tabular}
\caption{Model Uncertainty Analysis (Standard Deviations)}
\end{table}

\end{document}
